\section{Introduction} \label{sec:intro}
Computer vision and natural language understanding have made significant advances in recent years~\cite{kirillov2023segment,brown2020language}, where the availability of large datasets plays a critical role. Similarly, in robotics, scaling up human demonstration data has been key for learning new skills~\cite{brohan2022rt,padalkar2023open,fang2023rh20t}, with a clear trend of performance improvement with larger datasets~\cite{mandlekar2021matters,brohan2022rt}. However, human demonstrations, typically action-labeled trajectories collected via teleoperation devices~\cite{zhang2018deep,wu2023gello}, are time-consuming and labor-intensive to collect. For instance, collecting $130K$ trajectories in~\cite{brohan2022rt} took 17 months, making data collection a major bottleneck in robot learning.

Videos contain knowledge about behaviors, physics, and semantics, presenting an alternative data source. However, the lack of action labels makes utilization of video data in policy learning difficult. Previous works have addressed this by using self-supervised objectives for video pre-training to learn a feature representation of the observation for policy learning~\cite{sermanet2018tcn,nair2023r3m,seo2022reinforcement}. However, a feature representation only describes the state at the current time step, largely neglecting the transition dynamics that predicts future states. To explicitly model the transition dynamics, prior works have developed video prediction models that predict future image frames from current ones to guide policy learning~\cite{du2023unipi,yang2023learning, Escontrela23arXiv_VIPER}. However, learning a video prediction model for control introduces two challenges. Firstly, the task of video prediction avoids any abstraction by modeling changes to every pixel, coupling the physical motion with visual appearances such as texture, and lighting. This coupling makes modeling difficult, often resulting in hallucinations and unrealistic future predictions~\cite{du2023unipi}. Secondly, these models are computationally demanding in both training and inferencing. With limited computational resources, performance significantly declines. Moreover, the high inference cost compels these models to adopt open-loop execution~\cite{du2023unipi,black2023zero}, which tends to result in less robust policies.

In this paper, we propose a novel and structured representation to bridge video pre-training and policy learning. We first represent each state as a set of points in a video frame. To model the temporal structure in videos, we learn an Any-point Trajectory Model~(ATM) that takes the positions of the points in the current frame as input and outputs their future trajectories. We predict these trajectories in the camera coordinate frame to minimize the assumptions on calibrated cameras. These 2D point trajectories correspond to trajectories of particles in the 3D space and are a universal representation of the motions that can transfer to different domains and tasks. Contrasting with the video prediction approach of tracking changes in pixel intensity, the particle-based trajectory modeling offers a faithful abstraction of the physical dynamics, naturally incorporating inductive biases like object permanence and continuous motion. We first pre-train the trajectory model on actionless video datasets. After pre-training, the predicted trajectories can function as subgoals to guide the control policy. We then train trajectory-guided policies using only a minimal amount of action-labeled demonstration data. For training the ATMs, we generate self-supervised training data by leveraging recent advancements in vision models for accurate point tracking~\cite{karaev2023cotracker}. Across over \textbf{130} language-conditioned tasks we evaluated in both simulation and the real world, ATM significantly surpasses various strong baselines in video pre-training, achieving an average success rate of $63\%$ compared to the highest success rate of $37\%$ by previous methods, marking an improvement of over $80\%$. Additionally, we demonstrate that our method facilitates effective transfer learning from human videos and videos of a robot with a different morphology. We summarize our main contributions below:
\begin{enumerate}
  \item We propose an Any-point Trajectory Model, a simple and novel framework that bridges video pre-training to policy learning, leveraging the structured representation of particle trajectories.
  \item Through extensive experiments on simulated benchmarks and in the real world, we demonstrate that our method can effectively utilize video data in pre-training and significantly outperform various video pre-training baselines in an imitation learning setting.
  \item We demonstrate effective learning from cross-embodiment human and robot videos.
\end{enumerate}
